% Part: sets-functions-relations
% Chapter: sets
% Section: subsets

\documentclass[../../../include/open-logic-section]{subfiles}

\begin{document}

\olfileid[tr]{sfr}{set}{sub}
\olsection{Alt Kümeler ve Kuvvet Kümeleri}

\begin{explain}
Kümeleri sık sık karşılaştırmak isteyeceğiz. Doğal bir karşılaştırma şudur:
\emph{bir kümedeki her şey ötekinde de bulunur}. Bu durum,
yeni bir gösterim kullanmamızı gerektirecek kadar önemlidir.
\end{explain}

\begin{defn}[Alt Küme]
$A$ kümesinin her elemanı aynı zamanda $B$ kümesinin de elemanıysa, $A$
kümesine $B$ kümesinin bir \emph{alt kümesi} der ve $A \subseteq B$ yazarız.
$A$, $B$ kümesinin bir alt kümesi değilse $A \not\subseteq B$ yazarız.
$A \subseteq B$ ama $A \neq B$ ise $A \subsetneq B$ yazar ve $A$ kümesine
$B$ kümesinin bir \emph{öz alt kümesi} deriz.
\end{defn}

\begin{ex}
Her küme kendisinin bir alt kümesidir ve $\emptyset$ her kümenin bir
alt kümesidir. Çift doğal sayılar kümesi, doğal sayılar kümesinin bir
alt kümesidir. Ayrıca $\{ a, b \} \subseteq \{ a, b, c \}$ olur. Buna karşılık
$\{ a, b, e \}$, $\{ a, b, c \}$ kümesinin bir alt kümesi değildir.
\end{ex}

\begin{ex}
$2$ sayısı tam sayılar kümesinin bir elemanıyken çift sayılar kümesi, tam
sayılar kümesinin bir alt kümesidir. Bununla birlikte bir küme, başka bir
kümenin hem elemanı hem de alt kümesi olabilir.
Örneğin $\{0\} \in \{0, \{0\}\}$ ve aynı zamanda
$\{0\} \subseteq \{0, \{0\}\}$ olur.
\end{ex}

Kaplamsallık kümeler için bir özdeşlik ölçütü verir: $A = B$ eşitliği ancak ve
ancak $A$ kümesinin her elemanı $B$ kümesinin de elemanıysa ve bunun tersi de
geçerliyse sağlanır. ``Alt küme'' tanımı, $A \subseteq B$ ifadesini bu ölçütün
ilk yarısı olarak tanımlar: $A$ kümesinin her elemanı aynı zamanda $B$
kümesinin de elemanıdır. Tanım, $A$ ile $B$ kümelerinin yerleri
değiştirildiğinde de geçerlidir: $B \subseteq A$ ancak ve ancak $B$ kümesinin
her elemanı $A$ kümesinin de elemanıysa. Bu, kaplamsallıktaki ``tersine''
bölümüdür. Başka bir deyişle, kümeler ancak ve ancak birbirlerinin alt
kümeleriyse eşittir.

\begin{prop}
$A = B$ eşitliği ancak ve ancak hem $A \subseteq B$ hem de $B \subseteq A$
olduğunda geçerlidir.
\end{prop}

Şimdi birkaç kullanışlı gösterim daha tanıtabiliriz. $A$ kümesinin ne zaman
$B$ kümesinin bir alt kümesi olduğunu tanımlarken ``$A$ kümesinin her elemanı
\dots'' dedik ve ``$\dots$'' yerini ``$B$ kümesinin de elemanıdır''
ifadesiyle doldurduk. Bu anlatım \emph{biçimi} çok sık kullanıldığından onun
için biçimsel bir gösterim vermek yararlı olacaktır.

\begin{defn}\ollabel{forallxina}
$(\forall x \in A)\phi$, $\forall x(x \in A \lif \phi)$ ifadesinin
kısaltmasıdır. Benzer biçimde $(\exists x \in A)\phi$,
$\exists x(x \in A \land \phi)$ ifadesinin kısaltmasıdır.
\end{defn}

Bu gösterimle $A \subseteq B$ ifadesinin $(\forall x \in A)x \in B$
ifadesine denk olduğunu söyleyebiliriz.

Şimdi belirli türden bir kümeye geçiyoruz: verilmiş bir kümenin bütün
alt kümelerinden oluşan kümeye.

\begin{defn}[Kuvvet Kümesi]
Bir $A$ kümesinin bütün alt kümelerinden oluşan kümeye $A$ kümesinin
\emph{kuvvet kümesi} denir ve bu küme $\Pow{A}$ ile gösterilir.
  \[
    \Pow{A} = \Setabs{B}{B \subseteq A}
  \]
\end{defn}

\begin{ex}
$\{ a, b, c \}$ kümesinin olası bütün alt kümeleri nelerdir? Bunlar
$\emptyset$, $\{a \}$, $\{b\}$, $\{c\}$, $\{a, b\}$, $\{a, c\}$,
$\{b, c\}$ ve $\{a, b, c\}$ kümeleridir. Bu alt kümelerin tümünden oluşan
küme $\Pow{\{a,b,c\}}$'dir:
\[
\Pow{\{ a, b, c \}} = \{\emptyset, \{a \}, \{b\}, \{c\}, \{a, b\},
\{b, c\}, \{a, c\}, \{a, b, c\}\}
\]
\end{ex}

\begin{prob}
$\{a, b, c, d\}$ kümesinin bütün alt kümelerini listeleyin.
\end{prob}

\begin{prob}
$A$ kümesinin $n$ elemanı varsa $\Pow{A}$ kümesinin $2^n$ elemanı bulunduğunu
gösterin.
\end{prob}

\end{document}
